\documentclass[11pt]{article}

% ---- Portable preamble (compiles with XeLaTeX or pdfLaTeX, no project macros) ----
\usepackage{amsmath,amsthm,amssymb}
\usepackage[margin=1in]{geometry}
\usepackage[colorlinks=true,linkcolor=blue,urlcolor=blue,citecolor=blue]{hyperref}
\usepackage{mathtools}
\usepackage{enumitem}

\theoremstyle{plain}
\newtheorem{theorem}{Theorem}
\newtheorem{lemma}{Lemma}
\newtheorem{corollary}{Corollary}
\theoremstyle{definition}
\newtheorem{assumption}{Assumption}
\theoremstyle{remark}
\newtheorem{remark}{Remark}

\newcommand{\E}{\mathbb{E}}
\newcommand{\PP}{\mathbb{P}}
\newcommand{\Pn}{\mathbb{P}_n}
\newcommand{\R}{\mathbb{R}}
\newcommand{\cX}{\mathcal{X}}
\newcommand{\cT}{\mathcal{T}}
\newcommand{\cS}{\mathcal{S}}
\newcommand{\ez}{e_0}
\newcommand{\mz}{\mu_0}
\newcommand{\ehat}{\widehat{e}}
\newcommand{\mhat}{\widehat{\mu}}
\newcommand{\that}{\widehat{\tau}}
\newcommand{\thetahat}{\widehat{\theta}}
\newcommand{\Ltwo}[1]{\lVert #1 \rVert_{2}}
\DeclareMathOperator*{\argmin}{arg\,min}
\DeclarePairedDelimiter{\abs}{\lvert}{\rvert}
\DeclarePairedDelimiter{\card}{\lvert}{\rvert}

\title{Single-Tree ATT Inference on a Finite Covariate Space\\
\large The fixed finite regime}
\author{}
\date{2026-08-05}

\begin{document}
\maketitle

\begin{abstract}
We give valid $\sqrt n$ inference for the ATT when both nuisance functions are
estimated by a single globally-optimal decision tree fit on the \emph{full}
sample, with no sample splitting and no cross-fitting, on a covariate space that
is finite with fixed cardinality. The argument rests on one structural
observation: the induced estimand is \emph{invariant} across every partition fine
enough to represent the truth, so the data-adaptive choice of tree structure
cannot shift the target. Because the candidate set of trees is a fixed finite
set, uniformity over it is free, no entropy or tail conditions are needed, and
the ordinary empirical influence-function variance estimator is consistent. We
are explicit about what this does \emph{and does not} contribute
(Section~\ref{sec:contribution}): in this regime the inference problem is
finite-dimensional, and the content of the result is post-selection validity, not
nonparametric efficiency.
\end{abstract}

\section{Setup}\label{sec:setup}

Let $O=(X,A,Y)\sim P$ be i.i.d., $A\in\{0,1\}$, $Y\in\R$, and let the covariate
take values in a \emph{finite} product space
\[
  \cX \;=\; \prod_{j=1}^{d}\{1,\dots,K_j\},
  \qquad M \;:=\; \card{\cX} \;=\; \prod_j K_j .
\]
Throughout, $d$ and each $K_j$ are \textbf{fixed}: they do not grow with $n$.
Hence $M$ is a fixed finite constant. Write $\pi=\PP(A=1)$,
$\ez(x)=\PP(A=1\mid X=x)$, $\mz(x)=\E[Y\mid X=x,A=0]$, and let
\[
  \theta_0 \;=\; \E[Y(1)-Y(0)\mid A=1]
\]
be the ATT. The Neyman-orthogonal ATT score, in the empirical treated-mean
identification, is
\begin{equation}\label{eq:score}
  \psi(O;\theta,e,f)
  \;=\; \frac{A}{\pi}\bigl(Y-f(X)-\theta\bigr)
    \;-\; \frac{1}{\pi}\,\frac{e(X)\,(1-A)}{1-e(X)}\,\bigl(Y-f(X)\bigr).
\end{equation}
For candidate nuisances $(e,f)$ let $\theta(e,f)$ solve $\E_P[\psi]=0$. We use
throughout the standard bilinear remainder identity for doubly robust scores
(Robins--Rotnitzky; van der Laan--Robins; see also Kennedy's review), which for
\eqref{eq:score} reads
\begin{equation}\label{eq:bilinear}
  \theta(e,f)-\theta_0
  \;=\; \frac{1}{\pi}\,\E_P\!\left[\frac{\ez(X)-e(X)}{1-e(X)}\,
        \bigl(\mz(X)-f(X)\bigr)\right].
\end{equation}
Equation~\eqref{eq:bilinear} is not new and is not claimed as a contribution; it
is the textbook computation underlying double robustness, and we use it only as
bookkeeping.

\paragraph{Trees and partitions.}
A \emph{tree partition} $\tau$ of $\cX$ is a partition induced by recursive
axis-aligned binary splits on the coordinates of $\cX$. Because $\cX$ is finite
with fixed $M$, for any leaf budget $\bar L\le M$ the class
\[
  \cT_{\bar L} \;=\; \{\tau : \tau \text{ a tree partition of } \cX,\ \card{\tau}\le \bar L\}
\]
is a \textbf{fixed finite set}, and $\card{\cT_{\bar L}}$ is a constant not
depending on $n$. This single fact drives every simplification below.

For $\tau\in\cT_{\bar L}$ and $x\in\cX$ write $\ell_\tau(x)$ for the leaf
containing $x$. Given a partition $\tau$, the \emph{full-sample leaf-refit}
nuisances are
\begin{equation}\label{eq:refit}
  \ehat_\tau(x) = \frac{\sum_{i:\,X_i\in \ell_\tau(x)} A_i}{\#\{i: X_i\in\ell_\tau(x)\}},
  \qquad
  \mhat_\tau(x) = \frac{\sum_{i:\,X_i\in \ell_\tau(x),\,A_i=0} Y_i}
                        {\#\{i: X_i\in\ell_\tau(x),\,A_i=0\}} ,
\end{equation}
with $\ehat_\tau$ clipped to $[c,1-c]$ (Assumption~\ref{ass:causal}).

\paragraph{The estimator.}
Fix a budget $\bar L$ and a penalty $\lambda_n>0$. For each nuisance
$j\in\{e,\mu\}$ select
\begin{equation}\label{eq:select}
  \that_j \;\in\; \argmin_{\tau\in\cT_{\bar L}}
  \Bigl\{ R_n^{(j)}(\tau) + \lambda_n\card{\tau} \Bigr\},
\end{equation}
where $R_n^{(e)}$ is the empirical classification risk of the leaf-refit
propensity and $R_n^{(\mu)}$ the empirical squared-error risk of the leaf-refit
control-outcome model. Set $\ehat=\ehat_{\that_e}$, $\mhat=\mhat_{\that_\mu}$,
and let $\thetahat$ solve $\Pn[\psi(O;\theta,\ehat,\mhat)]=0$, i.e.\
\begin{equation}\label{eq:estimator}
  \thetahat
  = \frac{1}{n\widehat\pi}\sum_{i=1}^n\Bigl[
      A_i\bigl(Y_i-\mhat(X_i)\bigr)
      - \frac{\ehat(X_i)(1-A_i)}{1-\ehat(X_i)}\bigl(Y_i-\mhat(X_i)\bigr)\Bigr],
  \qquad \widehat\pi = \Pn[A].
\end{equation}
Note the \emph{minus} sign on the control-arm term; it follows from solving
$\Pn[\psi]=0$ with \eqref{eq:score} and must match the score.
\textbf{Both $\that_e$ and $\that_\mu$ are selected on all $n$ observations, and
both leaf refits use all $n$ observations. There is no sample splitting anywhere.}

\section{Assumptions}\label{sec:assumptions}

\begin{assumption}[Finite, fixed, fully supported covariate space]\label{ass:finite}
$\cX$ is finite with $M=\card{\cX}$ fixed in $n$, and $p_x:=\PP(X=x)>0$ for every
$x\in\cX$.
\end{assumption}

\begin{assumption}[Identification and strong positivity]\label{ass:causal}
Consistency; $\{Y(0),Y(1)\}\perp A\mid X$; and there is $c\in(0,\tfrac12]$ with
$\ez(x)\in[c,1-c]$ for every $x\in\cX$. Candidate propensities are clipped to
$[c,1-c]$.
\end{assumption}

\begin{assumption}[Second moments]\label{ass:moments}
$\E[Y^2\mid X=x]\le C_Y<\infty$ for every $x\in\cX$.
\end{assumption}

\begin{assumption}[Structural sparsity]\label{ass:sparsity}
For $j\in\{e,\mu\}$ let $\tau^*_j$ be a tree partition of minimal cardinality on
which $\gamma_{0,j}$ is constant ($\gamma_{0,e}=\ez$, $\gamma_{0,\mu}=\mz$), and
put $L^*_j=\card{\tau^*_j}$. The budget accommodates the truth:
$\bar L\ \ge\ \max(L^*_e,L^*_\mu)$.
\end{assumption}

\begin{assumption}[Global optimization]\label{ass:global}
$\that_j$ is a global minimizer of the penalized empirical risk over
$\cT_{\bar L}$, as in \eqref{eq:select}, with $\lambda_n\to 0$.
\end{assumption}

\begin{remark}[On Assumption~\ref{ass:sparsity}: this is sparsity, not correct
specification]\label{rem:not-exactness}
On a finite covariate space \emph{every} function is piecewise constant --- the
partition into singletons always works, so $L^*_j\le M$ holds automatically and
vacuously. Assumption~\ref{ass:sparsity} therefore asserts nothing about
functional form; its entire content is that $L^*_j$ is \emph{small enough to
display}, i.e.\ that the nuisances are structurally sparse. This should not be
described as an exactness or correct-specification assumption, and it is not
vulnerable to the objection ``you assumed the truth is a step function''.

Two further points. First, $L^*_j$ is the \emph{tree} complexity of
$\gamma_{0,j}$, which can be far smaller than $2^{s}$ for $s$ active
coordinates: dependence that is \emph{hierarchical} (the active coordinate set
differs by region) costs as few as $s+1$ leaves. Second, the assumption is
expensive precisely for \emph{additive} structure --- representing
$x_1+x_2+x_3$ exactly requires a leaf per configuration --- and cheap precisely
for region-varying active sets. That asymmetry is the substantive reason to use
trees rather than a sparse additive model, and it should be stated in the paper
with a concrete example rather than left as a verbal defence.
\end{remark}

\begin{remark}[On Assumption~\ref{ass:global}: global, not greedy]\label{rem:global}
Assumption~\ref{ass:global} is doing real work and cannot be weakened to greedy
recursive splitting. See Lemma~\ref{lem:margin} and
Remark~\ref{rem:greedy-no-margin}.
\end{remark}

\section{Sufficiency and invariance of the estimand}

\begin{lemma}[Sufficient class; invariance]\label{lem:invariance}
For $j\in\{e,\mu\}$ define the \emph{sufficient class}
\[
  \cS_j \;=\; \{\tau\in\cT_{\bar L} : \tau \text{ refines } \tau^*_j\}.
\]
Then:
\begin{enumerate}[nosep,label=(\roman*)]
  \item $\cS_j\neq\emptyset$; indeed $\tau^*_j\in\cS_j$ by
        Assumption~\ref{ass:sparsity}.
  \item For every $\tau\in\cS_j$ and every strictly positive weight function $w$,
        the $w$-weighted within-leaf population mean of $\gamma_{0,j}$ equals
        $\gamma_{0,j}$ pointwise.
  \item For every $(\tau_e,\tau_\mu)\in\cS_e\times\cS_\mu$, the induced
        working-model estimand satisfies
        $\theta(e_{\tau_e},\mu_{\tau_\mu})=\theta_0$ \emph{exactly}.
\end{enumerate}
\end{lemma}

\begin{proof}
(i) is immediate. For (ii), let $\ell$ be a leaf of $\tau\in\cS_j$. Since $\tau$
refines $\tau^*_j$, $\ell$ is contained in a single cell of $\tau^*_j$, on which
$\gamma_{0,j}$ is constant, say equal to $v$. Then for any $w>0$,
$\E[w\gamma_{0,j}\mid \ell]/\E[w\mid\ell] = v\,\E[w\mid\ell]/\E[w\mid\ell]=v$.
(The conditional expectations are well defined because $\PP(X\in\ell)>0$ by
Assumption~\ref{ass:finite}.) For (iii), by (ii) both population projections
equal the truth, so \emph{both} factors in \eqref{eq:bilinear} vanish; a fortiori
the product does. Note the two partitions need not coincide, and either factor
alone suffices.
\end{proof}

\begin{remark}[Why this is the load-bearing step]\label{rem:load-bearing}
Lemma~\ref{lem:invariance}(iii) says the target of inference does not move as the
selected structure varies over the sufficient class. Consequently we never need
$\that_j$ to recover $\tau^*_j$, and we never need a \emph{rate} for structure
selection: mere membership $\that_j\in\cS_j$ with probability tending to one
suffices. This is what removes the post-selection difficulty. It also disposes of
a distraction: $\mhat$ estimates the $(1-\ez)$-weighted within-leaf mean of
$\mz$, not the unweighted one, but by (ii) the two coincide on the sufficient
class, so the distinction is immaterial here.
\end{remark}

\section{The margin is free under global optimization}

\begin{lemma}[Positive margin]\label{lem:margin}
Let $R^{(j)}(\tau)$ denote the population risk of the best piecewise-constant
predictor on $\tau$ for nuisance $j$. Under
Assumptions~\ref{ass:finite}--\ref{ass:sparsity},
\[
  \Delta_j \;:=\; \min_{\tau\in\cT_{\bar L}\setminus\cS_j}
      \bigl\{R^{(j)}(\tau)-R^{(j)}(\tau^*_j)\bigr\} \;>\; 0 .
\]
\end{lemma}

\begin{proof}
Fix $\tau\in\cT_{\bar L}\setminus\cS_j$. Since $\tau$ does not refine $\tau^*_j$,
some leaf $\ell$ of $\tau$ contains $x,x'\in\cX$ with
$\gamma_{0,j}(x)\neq\gamma_{0,j}(x')$. Both have positive probability by
Assumption~\ref{ass:finite} (for $j=\mu$, positive probability in the control
arm, which holds because $p_x(1-\ez(x))\ge p_x c>0$ by
Assumption~\ref{ass:causal}). A single constant cannot equal two distinct values,
so the best constant on $\ell$ incurs strictly positive excess risk, whence
$R^{(j)}(\tau)-R^{(j)}(\tau^*_j)>0$. The minimum of finitely many strictly
positive numbers is strictly positive, and $\cT_{\bar L}$ is finite.
\end{proof}

\begin{remark}[Global optimization supplies the margin; greedy does
not]\label{rem:greedy-no-margin}
Lemma~\ref{lem:margin} is a statement about the \emph{global} risk functional,
and it is automatic --- not an assumption. This is a substantive advantage of
globally-optimal trees. Under greedy recursive splitting no such margin exists:
if $\ez$ depends on $X$ only through $\mathbf 1\{x_1=x_2\}$, then the marginal
risk reduction of \emph{either} correct first split is exactly zero, so a greedy
selector has no signal to follow and $\Delta=0$. The theorem below therefore
genuinely requires an exact solver, and this is one place where the paper's use
of optimal rather than CART-style trees is doing mathematical rather than
rhetorical work.
\end{remark}

\begin{lemma}[Selection lands in the sufficient class]\label{lem:selection}
Under Assumptions~\ref{ass:finite}--\ref{ass:global}, for $j\in\{e,\mu\}$
\[
  \PP\bigl(\that_j\in\cS_j\bigr)\;\longrightarrow\;1 .
\]
If in addition $Y$ is bounded, the convergence is exponential:
$\PP(\that_j\notin\cS_j)\le \card{\cT_{\bar L}}\,e^{-c_1 n\Delta_j^2}$ for a
constant $c_1>0$.
\end{lemma}

\begin{proof}
Differences of empirical risks across partitions are functions of the empirical
cell frequencies $\{\widehat p_x\}$ and the empirical cell-wise means of $A$ and
of $Y$ on the control arm; for squared error the $\Pn[Y^2]$ term is common to all
$\tau$ and cancels in any difference. By
Assumptions~\ref{ass:finite}--\ref{ass:moments} each of these finitely many
statistics converges in probability to its population counterpart (each cell and
each control cell has probability bounded below by a positive constant). Since
$\cT_{\bar L}$ is finite, convergence is automatically uniform:
\[
  \max_{\tau,\tau'\in\cT_{\bar L}}
  \bigl| \{R^{(j)}_n(\tau)-R^{(j)}_n(\tau')\} - \{R^{(j)}(\tau)-R^{(j)}(\tau')\}\bigr|
  \;=\; o_P(1).
\]
Also $\lambda_n\card{\tau}\le\lambda_n\bar L=o(1)$ uniformly. On the event that
the display above and the penalty are together smaller than $\Delta_j/2$, every
$\tau\notin\cS_j$ has strictly larger penalized empirical objective than
$\tau^*_j$ by Lemma~\ref{lem:margin}, so the minimizer lies in $\cS_j$. That
event has probability tending to one. If $Y$ is bounded, Hoeffding's inequality
applied to each of the finitely many statistics and a union bound over the fixed
finite class $\cT_{\bar L}$ give the stated exponential rate.
\end{proof}

\begin{remark}[Second moments suffice; no tail condition]\label{rem:no-tails}
Only $\PP(\that_j\in\cS_j)\to1$ is needed below, so
Assumption~\ref{ass:moments} suffices and no sub-Gaussian or bounded-outcome
condition is required. This is a consequence of $\cT_{\bar L}$ being a fixed
finite set: there is no growing entropy for the tail behaviour to have to beat.
\end{remark}

\begin{remark}[The penalty buys interpretability, not consistency]\label{rem:penalty}
By Lemma~\ref{lem:invariance} all members of $\cS_j$ induce the same estimand, so
the penalty in \eqref{eq:select} is statistically inert asymptotically: its role
is to select the \emph{smallest} sufficient tree, i.e.\ the most displayable one.
The regularization is doing interpretability work, not statistical work. This is
a pleasant alignment rather than a tension.
\end{remark}

\section{Finite-dimensional reduction and the main result}

\begin{lemma}[Conditional on a partition, the nuisance is
finite-dimensional]\label{lem:findim}
Fix $\tau\in\cT_{\bar L}$ and let $\eta_\tau=(\ehat_\tau,\mhat_\tau)$ as in
\eqref{eq:refit}. Then $\eta_\tau$ is determined by at most $2\bar L$ real
numbers (one leaf mean per nuisance per leaf), and under
Assumptions~\ref{ass:finite}--\ref{ass:moments}, if $\tau\in\cS_j$ for the
relevant $j$,
\[
  \max_{\ell\in\tau}\ \abs{\ehat_\tau(\ell)-\ez\!\restriction_\ell} = O_P(n^{-1/2}),
  \qquad
  \max_{\ell\in\tau}\ \abs{\mhat_\tau(\ell)-\mz\!\restriction_\ell} = O_P(n^{-1/2}),
\]
where $\gamma\!\restriction_\ell$ denotes the constant value of $\gamma$ on
$\ell$.
\end{lemma}

\begin{proof}
Each leaf is a union of cells, so $\PP(X\in\ell)\ge\min_x p_x>0$ and the leaf
count is $n\PP(X\in\ell)(1+o_P(1))\asymp n$; likewise the control count in
$\ell$ is $\asymp n$ by Assumption~\ref{ass:causal}. On $\cS_j$ the truth is
constant on $\ell$ (Lemma~\ref{lem:invariance}(ii)), so each leaf mean is an
average of $\asymp n$ i.i.d.\ variables with mean equal to that constant and
variance bounded by Assumption~\ref{ass:moments}; Chebyshev gives $O_P(n^{-1/2})$
per leaf, and $\card{\tau}\le\bar L$ is finite so the maximum has the same order.
\end{proof}

\begin{theorem}[Full-sample single-tree CLT]\label{thm:main}
Under Assumptions~\ref{ass:finite}--\ref{ass:global}, with $\thetahat$ as in
\eqref{eq:estimator},
\[
  \sqrt n\,\bigl(\thetahat-\theta_0\bigr)\;\rightsquigarrow\;N(0,V),
  \qquad
  V \;=\; \mathrm{Var}\bigl(\psi(O;\theta_0,\ez,\mz)\bigr),
\]
the semiparametric efficiency bound for the ATT. No sample splitting,
cross-fitting, or tail condition is used.
\end{theorem}

\begin{proof}
Let $E_n=\{\that_e\in\cS_e\}\cap\{\that_\mu\in\cS_\mu\}$; by
Lemma~\ref{lem:selection}, $\PP(E_n)\to1$. Since a sequence of random variables
that agrees with an asymptotically normal sequence on events of probability
tending to one has the same weak limit, it suffices to establish the expansion on
$E_n$.

Write $\psi(O;\theta,\eta) = \psi(O;\theta_0,\eta) - (A/\pi)(\theta-\theta_0)$
and recall $\E[A/\pi]=1$. Solving $\Pn[\psi(\cdot;\thetahat,\widehat\eta)]=0$ and
rearranging gives, with $\widehat\eta=(\ehat,\mhat)$,
\begin{equation}\label{eq:decomp}
  \thetahat-\theta_0
  = \frac{\pi}{\widehat\pi}\Bigl\{
    \underbrace{\Pn\bigl[\psi(O;\theta_0,\eta_0)\bigr]}_{\mathrm{(I)}}
    + \underbrace{(\Pn-P)\bigl[\psi(\cdot;\theta_0,\widehat\eta)-\psi(\cdot;\theta_0,\eta_0)\bigr]}_{\mathrm{(II)}}
    + \underbrace{P\bigl[\psi(\cdot;\theta_0,\widehat\eta)\bigr]}_{\mathrm{(III)}}
    \Bigr\}.
\end{equation}

\emph{Term (I).} An average of i.i.d.\ mean-zero terms with finite variance by
Assumptions~\ref{ass:causal}--\ref{ass:moments}; $\sqrt n\,\mathrm{(I)}
\rightsquigarrow N(0,V)$.

\emph{Term (III) is exactly the bilinear remainder.} Since $\theta(\ehat,\mhat)$
solves $P[\psi(\cdot;\theta,\ehat,\mhat)]=0$ and $\psi$ is affine in $\theta$
with $\E[A/\pi]=1$, we have $\mathrm{(III)}=\theta(\ehat,\mhat)-\theta_0$, so by
\eqref{eq:bilinear} and Assumption~\ref{ass:causal},
\[
  \abs{\mathrm{(III)}}
  \le \frac{1}{\pi c}\,\Ltwo{\ehat-\ez}\,\Ltwo{\mhat-\mz}.
\]
On $E_n$ both partitions are sufficient, so by Lemma~\ref{lem:findim} each factor
is $O_P(n^{-1/2})$ and $\mathrm{(III)}=O_P(n^{-1})=o_P(n^{-1/2})$. This is the
step at which structural sparsity pays: were the partitions not sufficient, each
factor would carry an irreducible approximation error and one would need the
product of those errors to be $o(n^{-1/2})$.

\emph{Term (II).} On $E_n$ condition on the selected pair
$(\that_e,\that_\mu)\in\cS_e\times\cS_\mu$. Because $\cS_e\times\cS_\mu$ is a
\textbf{fixed finite} set, it suffices to bound the term for each fixed pair and
take a maximum over finitely many bounds; a maximum of finitely many $o_P(n^{-1/2})$
quantities is $o_P(n^{-1/2})$. \emph{No entropy integral, covering number, or
union bound over a growing class is required.} Fix such a pair. By
Lemma~\ref{lem:findim} the nuisance is a vector
$\widehat v\in\R^{2\bar L}$ of leaf means with
$\widehat v - v_0 = O_P(n^{-1/2})$, and on the clipped region $[c,1-c]$ the map
$v\mapsto\psi(\cdot;\theta_0,\eta_v)$ is affine in the outcome-leaf coordinates
and smooth (with derivatives bounded by $c^{-2}$ and Assumption~\ref{ass:moments})
in the propensity-leaf coordinates. A first-order expansion gives
\[
  \psi(\cdot;\theta_0,\eta_{\widehat v})-\psi(\cdot;\theta_0,\eta_{v_0})
  = \sum_{k=1}^{2\bar L} (\widehat v_k - v_{0,k})\, g_k(\cdot) \;+\; r(\cdot),
  \qquad \Ltwo{r} = O_P\bigl(\abs{\widehat v - v_0}^2\bigr),
\]
for \emph{fixed} functions $g_k$ depending only on the (now fixed) partitions.
Hence
\[
  \mathrm{(II)}
  = \sum_{k=1}^{2\bar L} (\widehat v_k - v_{0,k})\,(\Pn-P)[g_k]
    + (\Pn-P)[r]
  = O_P(n^{-1/2})\cdot O_P(n^{-1/2}) + o_P(n^{-1/2})
  = o_P(n^{-1/2}),
\]
since each $(\Pn-P)[g_k]=O_P(n^{-1/2})$ by the central limit theorem for the
fixed square-integrable function $g_k$, and there are finitely many $k$.

\emph{Conclusion.} $\widehat\pi-\pi=O_P(n^{-1/2})$ and
$\mathrm{(I)}=O_P(n^{-1/2})$, so $(\pi/\widehat\pi-1)\mathrm{(I)}=O_P(n^{-1})$.
Combining, $\sqrt n(\thetahat-\theta_0)=\sqrt n\,\mathrm{(I)}+o_P(1)
\rightsquigarrow N(0,V)$.
\end{proof}

\begin{corollary}[Variance estimation and confidence intervals]\label{cor:variance}
Under the conditions of Theorem~\ref{thm:main}, the ordinary empirical
influence-function variance estimator
\[
  \widehat V \;=\; \Pn\bigl[\psi(O;\thetahat,\ehat,\mhat)^2\bigr]
\]
satisfies $\widehat V\to_P V$, and
$\thetahat\pm z_{1-\alpha/2}\sqrt{\widehat V/n}$ is an asymptotically valid
$(1-\alpha)$ confidence interval for $\theta_0$.
\end{corollary}

\begin{proof}
On $E_n$, $\widehat v\to_P v_0$ and $\thetahat\to_P\theta_0$; on the clipped
region $\psi^2$ is a continuous function of $(\theta,v)$ and of $O$, dominated
under Assumptions~\ref{ass:causal}--\ref{ass:moments}. Since
$\cS_e\times\cS_\mu$ is finite, the convergence is uniform over it, and
$\PP(E_n)\to1$. Apply the law of large numbers and continuity.
\end{proof}

\begin{remark}[This resolves the missing variance estimator]\label{rem:variance-resolved}
No leave-one-out or jackknife construction is needed in this regime: the plain
plug-in influence-function variance is consistent, and it is already what the
software computes. The delicate leaf-wise conditional-variance argument required
for the cross term in the general (growing) regime is likewise unnecessary here,
because the crude finite-dimensional bound in Term (II) already suffices.
\end{remark}

\begin{remark}[A falsifiable prediction for the simulations]\label{rem:prediction}
Corollary~\ref{cor:variance} asserts that the naive Wald interval for the
single-tree estimator is asymptotically valid. Earlier internal experiments found
it \emph{under}-covering. Theorem~\ref{thm:main} attributes that entirely to
finite-sample selection instability, i.e.\ to $\PP(E_n)$ being appreciably below
one at moderate $n$, and therefore predicts that undercoverage should shrink as
the measured rate of ``selected partition is sufficient'' approaches one. That
rate is now instrumented directly. If undercoverage persists while sufficiency
approaches one, the theorem is wrong or an assumption fails --- most plausibly
Assumption~\ref{ass:global}, if the deployed solver is not attaining the global
optimum. This is a genuine test rather than a consistency check.
\end{remark}

\section{What this does and does not contribute}\label{sec:contribution}

We state the limitation squarely, because a referee will.

\begin{remark}[The inference problem here is finite-dimensional]\label{rem:findim-objection}
Under Assumption~\ref{ass:finite} with $M$ fixed, $P$ is determined by finitely
many parameters and the ATT is a smooth function of them. Theorem~\ref{thm:main}
is therefore \emph{not} a nonparametric efficiency result, and it should not be
sold as one. Conditional on a sufficient partition the estimator is an ordinary
smooth $Z$-estimator in $\le 2\bar L+1$ parameters --- which is precisely why no
sample splitting is required. Claiming otherwise would invite, correctly, the
objection that cross-fitting was never needed in the first place.
\end{remark}

What the theorem does contribute, stated so as to survive
Remark~\ref{rem:findim-objection}:
\begin{enumerate}[label=(C\arabic*),leftmargin=*]
  \item \textbf{Post-selection validity for tree structure.} The structures
  $\that_e,\that_\mu$ are chosen by data-adaptive optimization on the same sample
  used for inference. Naive post-selection inference is not generally valid;
  here it is, and the reason is specific and checkable ---
  Lemma~\ref{lem:invariance}: the estimand is invariant over the sufficient
  class, so selection cannot move the target. The existing post-selection
  inference literature is built for linear-model variable selection, not for
  tree-structure selection, so the analogue does not appear to exist.
  \item \textbf{The displayed model is the estimator.} Not a post hoc surrogate
  for a black-box nuisance, and not $K$ different trees. One propensity tree and
  one outcome tree, fit on every observation, are simultaneously the audit object
  and the inferential object.
  \item \textbf{Global optimization is necessary, and we can say why.}
  Remark~\ref{rem:greedy-no-margin}: greedy selection has zero margin under
  interaction-dominated truths, so the guarantee genuinely depends on using an
  exact solver.
  \item \textbf{Regularization serves interpretability at no statistical cost}
  (Remark~\ref{rem:penalty}).
\end{enumerate}

\begin{remark}[Where the nonparametric content lives]\label{rem:where-np}
The nonparametric content of the paper sits in the cross-fitted result for
optimal trees on continuous covariates (the companion section), where tree
complexity must grow with $n$ and approximation error is irreducible. The present
theorem is the interpretable, finite-support counterpart. Readers who do not need
a single displayed tree should use the cross-fitted estimator; that is the honest
answer to ``how would you relax Assumption~\ref{ass:sparsity}''.
\end{remark}

\section{Open items}

\begin{enumerate}[nosep,label=(\arabic*)]
  \item \textbf{Growing regime.} If $M=M_n\to\infty$ (many-valued or continuous
  coordinates) or $\bar L=\bar L_n\to\infty$, then $\cT_{\bar L}$ is no longer
  fixed, uniformity is no longer free, and the entropy $\bar L_n\log(dK_n)$
  enters, along with a tail condition and possibly sample splitting for structure
  selection. Deliberately not treated here.
  \item \textbf{Graded case} $L^*_j>\bar L$. Then
  Lemma~\ref{lem:invariance} fails, Term (III) no longer vanishes, and one needs
  $\delta_e\delta_\mu=o(n^{-1/2})$ where $\delta_j$ is the approximation error of
  the best $\bar L$-leaf tree. On a finite space no decay rate is implied by the
  domain (parity on $d$ binary coordinates has $\delta(L)$ bounded away from zero
  until $L\ge2^{d-1}$), so this requires a separate assumption. Note that here
  $\delta_j$ is \emph{estimable} whenever $M\ll n$, via the excess risk of the
  $\bar L$-leaf fit relative to the saturated cell-wise fit, which permits a
  bias-aware interval instead of an assumed rate.
  \item \textbf{Assumption~\ref{ass:global} versus the deployed solver.} Depth
  caps and solver timeouts may prevent attainment of the global optimum over
  $\cT_{\bar L}$. This is an empirical question, currently open, and it is the
  most likely failure point of the theorem in practice.
\end{enumerate}

\section{Nuisances outside the class: honest inference at a stated
price}\label{sec:outside}

Assumption~\ref{ass:finite} fails in three ways of practical interest: the number
of covariates grows ($d=d_n\to\infty$), the number of levels grows
($K_j=K_{j,n}\to\infty$), or some covariates are continuous ($M=\infty$).
Assumption~\ref{ass:sparsity} fails whenever $L^*_j>\bar L$, i.e.\ whenever the
truth is a larger tree than we are willing to display. This section states what
survives. The short answer is that \emph{validity} survives and \emph{efficiency}
does not, and that the price is estimable.

\subsection{Two remainders fail, with different characters}

It is worth separating them, because they call for different treatment.

\begin{itemize}[leftmargin=*]
  \item \textbf{Term (III), the bias, is a property of the truth.} Outside the
  class $\delta_e\delta_\mu\neq0$, where
  $\delta_j := \inf_{\tau\in\cT_{\bar L}}\Ltwo{\gamma_{\tau,j}-\gamma_{0,j}}$ is
  the approximation error of the best admissible tree ($\delta_\mu$ defined on the
  $(1-\ez)$-weighted projection, which by Assumption~\ref{ass:causal} is
  norm-equivalent to the unweighted one up to $1/c$). Nothing we do
  algorithmically changes $\delta_j$. It must be assumed or estimated.
  \item \textbf{Term (II), the empirical process, is enforceable by the
  algorithm.} In the fixed finite regime it was free for three reasons, all of
  which now fail: the candidate class was fixed and finite (so uniformity cost
  nothing), leaf probabilities were bounded below by a constant, and the nuisance
  had fixed dimension. Outside the class the entropy $\bar L_n\log(dK_n)$ enters,
  a tail condition is needed for the associated exponential inequality, and
  $\E\Ltwo{\mhat-\mz}^2 \asymp n^{-1}\sum_\ell \sigma^2_\ell/p_\ell$ is
  $\asymp \bar L_n/n$ \emph{only} if $\min_\ell p_\ell\gtrsim 1/\bar L_n$.
\end{itemize}

Accordingly, bounding the bias is \textbf{necessary but not sufficient}. The
following are conditions we impose on the estimator, not on $P$, and they must be
imposed rather than assumed away: a leaf budget $\bar L_n = o(\sqrt n)$; a minimum
leaf count $m_n$ with $\bar L_n m_n\lesssim n$; and clipping
$\ehat\in[c,1-c]$. None is currently enforced or measured in the implementation.

\subsection{Can the bias be bounded? Three regimes}

\begin{enumerate}[leftmargin=*,label=(\alph*)]
  \item \textbf{Growing $d$ or $K$, still discrete, $M_n\ll n$.} Then $\delta_j$
  is \emph{estimable}: it is the excess risk of the $\bar L$-leaf fit relative to
  the saturated cell-wise fit, available at rate $O(\sqrt{M_n/n})$. No assumption
  about the truth is required. This is the favourable case.
  \item \textbf{Continuous covariates.} The saturated reference does not exist and
  $\delta_j$ is \textbf{not estimable from the data alone}. One can measure the
  approximation error visible at whatever resolution is fittable, but never
  exclude structure below it. This is the classical obstruction in the honest
  confidence interval literature and it is why bias-aware approaches
  (Armstrong--Koles\'ar) assume a smoothness constant, which is not testable.
  \item \textbf{Any regime, using a $\sqrt n$-consistent anchor.} If a
  $\sqrt n$-consistent estimator of $\theta_0$ is available --- and the
  cross-fitted optimal-tree estimator is exactly that, under its own conditions
  --- then the single tree's bias is estimable as a difference, with no smoothness
  assumption. This is Proposition~\ref{prop:honest} and it is the route we
  recommend, because the anchor is already part of the paper.
\end{enumerate}

\subsection{Honest inference via an anchor}

\begin{lemma}[Do not subtract; widen]\label{lem:no-subtract}
Let $\widetilde\theta$ be the cross-fitted estimator and
$\widehat\delta = \thetahat-\widetilde\theta$. Then
$\thetahat-\widehat\delta = \widetilde\theta$ identically. Hence any procedure
that \emph{subtracts} $\widehat\delta$ from the displayed-tree estimate returns
the cross-fitted estimator exactly, discarding the interpretable point estimate
and gaining nothing.
\end{lemma}

Lemma~\ref{lem:no-subtract} is trivial but it settles a question that has caused
confusion: $\widehat\delta$ must be used to \emph{widen the interval}, never to
correct the point estimate. The point estimate remains the displayed tree; the
interval pays for the gap.

\begin{theorem}[Honest interval for $\theta_0$ outside the class]\label{prop:honest}
Let $\widetilde\theta$ satisfy
$\sqrt n(\widetilde\theta-\theta_0)\rightsquigarrow N(0,\widetilde V)$ with
$\widetilde\sigma^2\to_P\widetilde V$, and let
$\widehat\delta=\thetahat-\widetilde\theta$. Define
\[
  C_n \;=\; \Bigl[\;\thetahat \;\pm\; \bigl(\,\abs{\widehat\delta}
      + z_{1-\alpha/2}\,\widetilde\sigma/\sqrt n\,\bigr)\Bigr].
\]
Then $\liminf_n \PP(\theta_0\in C_n)\ \ge\ 1-\alpha$, with no condition on
$\delta_e,\delta_\mu$, and in particular with no smoothness assumption on the
nuisances.
\end{theorem}

\begin{proof}
Write $\thetahat-\theta_0 = \widehat\delta + (\widetilde\theta-\theta_0)$, so
$\abs{\thetahat-\theta_0}\le\abs{\widehat\delta}+\abs{\widetilde\theta-\theta_0}$
by the triangle inequality. Hence on the event
$\{\abs{\widetilde\theta-\theta_0}\le z_{1-\alpha/2}\widetilde\sigma/\sqrt n\}$ we
have $\abs{\thetahat-\theta_0}\le\abs{\widehat\delta}+z_{1-\alpha/2}\widetilde\sigma/\sqrt n$,
i.e.\ $\theta_0\in C_n$. That event has probability tending to $1-\alpha$ by the
anchor's asymptotic normality and the consistency of $\widetilde\sigma$.
\end{proof}

\begin{corollary}[The price, and where it plateaus]\label{cor:width}
The width of $C_n$ is $2(\abs{\widehat\delta}+z_{1-\alpha/2}\widetilde\sigma/\sqrt n)$,
and $\abs{\widehat\delta} = O_P\bigl(\delta_e\delta_\mu + n^{-1/2}\bigr)$ by
\eqref{eq:bilinear} applied to the tree nuisances. Therefore
\[
  \mathrm{width}(C_n) \;\asymp_P\; \max\bigl\{\,n^{-1/2},\ \delta_e\delta_\mu\,\bigr\}.
\]
The interval contracts at the parametric rate \emph{if and only if}
$\delta_e\delta_\mu = O(n^{-1/2})$; otherwise its width \textbf{plateaus} at the
bias floor.
\end{corollary}

\begin{remark}[This is the interpretability--accuracy trade-off, quantified]
Corollary~\ref{cor:width} is the result we want, and it replaces an earlier and
unsuccessful attempt to prove an impossibility theorem. The correct statement is
not ``a bounded-complexity interpretable tree cannot deliver $\sqrt n$ inference
for $\theta_0$'' --- which requires a minimax lower bound that was never
established, and which is false in the regime of Theorem~\ref{thm:main} --- but
rather: \emph{$\theta_0$ inference remains valid, and the cost of insisting on a
displayable tree is an interval whose width contracts at $\sqrt n$ only while the
tree is rich enough to track the nuisances, and plateaus thereafter.} The
plateau height, $\delta_e\delta_\mu$, is the price of interpretability, and by
\S\ref{sec:outside}(a) or (c) it is estimable rather than merely assumed.
\end{remark}

\begin{remark}[Conservativeness, and the sharper construction]\label{rem:conservative}
Theorem~\ref{prop:honest} is deliberately crude: $C_n$ contains the anchor's own
interval recentred at $\thetahat$, so it is never narrower than the cross-fitted
interval. The gain is the interpretable point estimate, not efficiency. A sharper
bias-aware critical value (Armstrong--Koles\'ar) exploits the fact that a bias of
known magnitude $B$ costs less than $B$ in interval half-width, and should be used
in practice; Theorem~\ref{prop:honest} is stated in the crude form because it
requires nothing beyond the triangle inequality.
\end{remark}

\begin{remark}[Implementation gap: the sampling error of $\widehat\delta$ is
currently ignored]\label{rem:se-delta}
The shipped implementation sets $\mathrm{se}(\widehat\delta)=0$, on the stated
grounds that this yields ``the tightest interval consistent with the guarantee''.
That reasoning holds only if $\widehat\delta$ is a \emph{known} bound. It is an
estimate, with sampling error of order $n^{-1/2}$, so treating it as exact risks
undercoverage. Theorem~\ref{prop:honest} avoids the issue by adding the anchor's
full $z\widetilde\sigma/\sqrt n$ on top of $\abs{\widehat\delta}$; any tighter
construction must account for $\mathrm{se}(\widehat\delta)$ explicitly.
\end{remark}

\begin{remark}[Preliminary empirical support, with a caveat]\label{rem:empirics-outside}
On the existing \texttt{complex}-DGP runs (355 replications, $n$ from 500 to
10{,}000) the bias-aware interval dominates the naive Wald interval at every $n$
and repairs the worst cell: coverage of $\theta_0$ was $0.762$ (naive) versus
$0.917$ (bias-aware) at $n=5000$, and was never worse elsewhere. Two caveats. The
comparison is between two interval constructions computed from the \emph{same}
point estimates, which makes it relatively robust; but those point estimates were
produced before a leaf-value defect in the tree-refit routine was corrected, and
coverage was computed on converged replications only, which is not
missing-at-random. The comparison should be regenerated. Independently of
coverage, the same runs show $\sqrt n\times\mathrm{bias}$ \emph{flat} in $n$ for
the single-tree estimator ($0.32$--$0.59$ across a twentyfold increase in $n$)
while collapsing for the cross-fitted estimator ($0.30\to0.035$) --- the
signature of $\delta_e\delta_\mu$ failing to be $o(n^{-1/2})$, i.e.\ of
Assumption~\ref{ass:sparsity} not holding for the partitions the solver actually
returned. Note that coverage is a comparatively insensitive instrument for this:
with bias-to-half-width ratios of $0.17$--$0.31$ the coverage loss is only a few
points, whereas $\sqrt n\times\mathrm{bias}$ shows the failure directly.
\end{remark}

\subsection{Three tiers of claim}

We recommend presenting exactly three, and being explicit that they answer
different questions.

\begin{enumerate}[leftmargin=*,label=\textbf{T\arabic*.}]
  \item \textbf{Inside the class} (Assumptions~\ref{ass:finite}--\ref{ass:global}):
  efficient $\sqrt n$ inference for $\theta_0$ from the displayed trees, ordinary
  plug-in variance, no splitting (Theorem~\ref{thm:main},
  Corollary~\ref{cor:variance}).
  \item \textbf{Outside the class, bias bounded}: valid but inefficient $\theta_0$
  inference, with width $\asymp\max\{n^{-1/2},\delta_e\delta_\mu\}$
  (Theorem~\ref{prop:honest}, Corollary~\ref{cor:width}). The displayed trees
  remain the point estimate.
  \item \textbf{Outside the class, bias not bounded}: report the displayed trees
  as an audit object alongside the cross-fitted estimate, with $\widehat\delta$ as
  a fidelity diagnostic and \emph{no} inferential claim attached to the trees.
  This is the honest fallback and it should be stated as available, so that a
  practitioner who cannot certify (a) or (c) of \S\ref{sec:outside} is not left
  without a defensible option.
\end{enumerate}

\end{document}
